\documentclass{scrbook}
\title{The Span Gauge}
\author{Paper Bridge Workshop}
\date{}
\begin{document}
\frontmatter
\maketitle
\tableofcontents
\mainmatter
\part{Making the gauge}
\chapter{The sliding strip}
\section{Cutting the slot}
Mark the centre of the card before cutting the slot.\footnote{Keep the offcut as a shim.}
\subsection{Checking the edge}
Use the dark strip in Figure~\ref{fig:shim} to check for daylight.
\begin{figure}[htbp]
\centering\rule{28mm}{3mm}
\caption{The shim seen on edge}\label{fig:shim}
\end{figure}
\begin{itemize}\item Cut across the grain.\item Keep the straight edge.\end{itemize}
\newpage
\minisec{Bench note}
Store the strip flat until the glue has set.
\chapter{Setting the stops}
\section{The second mark}
\subsection{Trial positions}
Table~\ref{tab:stops} records the positions. Equation~\ref{eq:gap} gives the gap.
\begin{equation}\label{eq:gap}
g=b-a
\end{equation}
\begin{table}[htbp]
\centering\caption{Stop positions}\label{tab:stops}
\begin{tabular}{lr}Stop & Position (mm)\\Near & 6\\Far & 18\end{tabular}
\end{table}
\begin{enumerate}\item Seat the first stop.\item Slide the strip to the second mark.\end{enumerate}
\appendix
\chapter{Spare gauges}
Keep a spare strip with each bridge.
\backmatter
\addchap{Workshop record}
The folded gauge follows the bench card~\cite{card}.
\begin{thebibliography}{9}
\bibitem{card} Paper Bridge Workshop, \emph{Stop setting card}.
\end{thebibliography}
\end{document}
